% =====================================================================
%  Cadenza — Documento principal de tesis (versión fusionada)
%  Plantilla: IEEEtran (conference)
%
%  Contenido actual (única fuente de verdad — no editar introduccion.tex):
%    I.   Introducción
%    II.  Planteamiento del Problema
%
%  Secciones pendientes de fusionar aquí cuando existan:
%    III. Revisión sistemática de literatura (docs/revision-literatura-prisma.md)
%    IV.  Arquitectura general DBB          (docs/arquitectura-dbb.md)
%
%  Bibliografía: UNA SOLA lista al final (13 entradas); todas las citas
%  del texto apuntan a estas claves — sin referencias huérfanas.
%
%  Compilación (pdflatex, sin bibtex — bibliografía inline):
%    pdflatex main
%
%  Referencias VERIFICAR (confirmar autores/DOI antes de entrega):
%    riosvila2026, mlsp2025, ricordi2024  (ver docs/literatura/)
% =====================================================================
\documentclass[conference]{IEEEtran}
\IEEEoverridecommandlockouts

\usepackage{cite}
\usepackage{amsmath,amssymb,amsfonts}
\usepackage{graphicx}
\usepackage{textcomp}
\usepackage{xcolor}
\usepackage{array}
% Paquetes para español:
\usepackage[utf8]{inputenc}
\usepackage[T1]{fontenc}
\usepackage[spanish,es-tabla]{babel}
% Diagramas (árbol del problema):
\usepackage{tikz}
\usetikzlibrary{positioning,arrows.meta}

\begin{document}

% ---------------------------------------------------------------------
% TÍTULO
% ---------------------------------------------------------------------
\title{Cadenza: reconocimiento, validación y aprendizaje en la
digitalización de partituras}

% ---------------------------------------------------------------------
% AUTORES
% ---------------------------------------------------------------------
\author{\IEEEauthorblockN{Andr\'es Carrero Mart\'inez}
\IEEEauthorblockA{\textit{Ingenier\'ia de Sistemas} \\
\textit{Universidad del Norte}\\
Barranquilla, Colombia \\
aacarrero@uninorte.edu.co}
\and
\IEEEauthorblockN{Sebasti\'an Ib\'a\~nez R\'ios}
\IEEEauthorblockA{\textit{Ingenier\'ia de Sistemas} \\
\textit{Universidad del Norte}\\
Barranquilla, Colombia \\
sdibanez@uninorte.edu.co}
\and
\IEEEauthorblockN{Carlos Herrera Rojano}
\IEEEauthorblockA{\textit{Ingenier\'ia de Sistemas} \\
\textit{Universidad del Norte}\\
Barranquilla, Colombia \\
carlosrojano@uninorte.edu.co}
}

\maketitle

% ---------------------------------------------------------------------
% FIGURA 1: árbol del problema (flotante a página completa).
% Se declara aquí arriba para que LaTeX la ubique al inicio de las
% primeras páginas y NO caiga sobre la bibliografía al final.
% ---------------------------------------------------------------------
\begin{figure*}[!t]
\centering
\begin{tikzpicture}[
  font=\footnotesize,
  node distance=0.55cm and 1.0cm,
  caja/.style={draw, thick, align=center, text width=4.6cm, inner sep=5pt},
  causa/.style={caja, fill=red!12},
  problema/.style={caja, fill=orange!30, text width=4.4cm, font=\footnotesize\bfseries},
  efecto/.style={caja, fill=blue!10}
]
\node[causa] (c1) {Salida OMR ruidosa: alturas, duraciones, armaduras y s\'imbolos err\'oneos};
\node[causa, below=of c1] (c2) {Las correcciones del usuario se descartan: el sistema no mejora con el uso};
\node[causa, below=of c2] (c3) {Ausencia de verificaci\'on autom\'atica de reglas de teor\'ia musical};
\node[problema, right=of c2] (p) {Digitalizaci\'on de partituras lenta y poco confiable: correcci\'on manual s\'imbolo a s\'imbolo en editores convencionales};
\node[efecto, right=of p] (e2) {Transcripciones MusicXML con inconsistencias estructurales};
\node[efecto, above=of e2] (e1) {Cuellos de botella en la preservaci\'on y difusi\'on de archivos musicales};
\node[efecto, below=of e2] (e3) {Abandono de herramientas OMR y retorno a la transcripci\'on manual};
\draw[-{Stealth}, thick] (c1.east) -- (p.west);
\draw[-{Stealth}, thick] (c2.east) -- (p.west);
\draw[-{Stealth}, thick] (c3.east) -- (p.west);
\draw[-{Stealth}, thick] (p.east) -- (e1.west);
\draw[-{Stealth}, thick] (p.east) -- (e2.west);
\draw[-{Stealth}, thick] (p.east) -- (e3.west);
\end{tikzpicture}
\caption{\'Arbol del problema de la digitalizaci\'on asistida de partituras.}
\label{fig:arbol}
\end{figure*}

% =====================================================================
% I. INTRODUCCIÓN
% =====================================================================
\section{Introducci\'on}

La digitalizaci\'on del patrimonio musical constituye una necesidad
creciente para archivos, bibliotecas, academias y agrupaciones musicales
que conservan partituras f\'isicas en papel. El reconocimiento \'optico de
m\'usica (OMR, por sus siglas en ingl\'es) es la disciplina que aborda la
conversi\'on autom\'atica de im\'agenes de partituras en representaciones
simb\'olicas editables, y ha sido objeto de investigaci\'on durante m\'as de
cinco d\'ecadas \cite{omr2020}, evolucionando desde sistemas basados en
reglas hasta \emph{pipelines} de aprendizaje profundo que transcriben
partituras completas de forma autom\'atica
\cite{rebelo2012}\cite{calvo2018}\cite{riosvila2026}.

A pesar de estos avances, los sistemas OMR modernos cometen errores que
impiden su uso directo en entornos reales: notas con alturas o duraciones
incorrectas, compases desbalanceados, armaduras inconsistentes y s\'imbolos
omitidos \cite{calvo2018}. Los modelos de reconocimiento se estudian
predominantemente desde la precisi\'on de su salida autom\'atica, mientras
que la integraci\'on del reconocimiento en el trabajo pr\'actico de un
transcriptor ha recibido mucha menos atenci\'on \cite{rizo2023}. Como
consecuencia, quien digitaliza partituras termina revisando manualmente
cada s\'imbolo en un editor musical convencional, un proceso tan lento como
la transcripci\'on manual, o acepta resultados err\'oneos sin verificarlos.
En el caso de los manuscritos musicales modernos, esta situaci\'on se
agudiza, pues la variabilidad de la escritura degrada a\'un m\'as la
confiabilidad del reconocimiento autom\'atico \cite{ricordi2024}. Del mismo
modo, el aprovechamiento de las correcciones del usuario para mejorar el
sistema mediante ciclos de aprendizaje activo sigue siendo poco explorado
en OMR, donde los primeros estudios arrojan resultados incluso
contradictorios \cite{mlsp2025}.

% COMPLETAR: párrafo de contexto local — personalizar con datos reales de la
% institución/colección objetivo (ver docs/introduccion-borrador.md, bloque 3).
En el contexto global, instituciones conservan colecciones de partituras
impresas y escritas a mano que requieren digitalizaci\'on para su
preservaci\'on, consulta y difusi\'on. La transcripci\'on manual de estas
colecciones representa un costo de tiempo considerable, y las herramientas
de reconocimiento existentes no ofrecen un flujo de trabajo que garantice
la calidad del resultado. La digitalizaci\'on de colecciones musicales ha
sido abordada recientemente desde la arch\'istica con metodolog\'ias de
anotaci\'on y clasificaci\'on autom\'atica \cite{ricordi2024}; a\'un as\'i,
faltan herramientas de transcripci\'on asistida pensadas para el trabajo
del transcriptor.

Para la evaluaci\'on del sistema se dispone de conjuntos de datos
p\'ublicos de referencia que cubren el alcance del proyecto. PrIMuS y
Camera-PrIMuS \cite{calvo2018} proveen 87\,678 incipits de partituras
impresas monof\'onicas, el segundo con distorsiones que simulan fotograf\'ias
tomadas en condiciones reales, y \emph{ground truth} simb\'olico en
m\'ultiples formatos. El Sheet Music Benchmark (SMB)
\cite{smb2026} estandariza la evaluaci\'on con 685 p\'aginas que incluyen
monofon\'ia y \emph{pianoform}, divisiones predefinidas y una m\'etrica de
error detallada por elemento musical (OMR-NED). Para el subconjunto
manuscrito, MUSCIMA++ \cite{muscima2017} ofrece m\'as de 91\,000 s\'imbolos
anotados sobre partituras escritas a mano por m\'usicos actuales. La
disponibilidad de \emph{ground truth} anotado en estos conjuntos permite
evaluar la plataforma de forma cuantitativa y reproducible, siguiendo los
protocolos de la literatura de OMR. De manera complementaria, se contempla
un piloto sobre un conjunto reducido de partituras reales (fotograf\'ias)
para evaluar la usabilidad del flujo completo de correcci\'on, sin requerir
\emph{ground truth} anotado.

Desde el punto de vista disciplinar, un sistema de digitalizaci\'on
asistida de partituras puede estructurarse sobre tres pilares: el
\emph{pipeline} cl\'asico de OMR (preprocesado, detecci\'on de pentagramas,
detecci\'on y clasificaci\'on de s\'imbolos, ensamblaje y reconstrucci\'on
musical \cite{omr2020}), la validaci\'on de la salida
mediante reglas de teor\'ia musical \cite{cuthbert2010}, y la correcci\'on
asistida por un humano en el ciclo \cite{amershi2014}. La validaci\'on
autom\'atica permite detectar errores estructurales (m\'etrica, duraciones,
armaduras) sin intervenci\'on humana, reduciendo el espacio de revisi\'on a
los puntos sospechosos. El aprendizaje activo complementa el ciclo al
seleccionar qu\'e correcciones del usuario tienen mayor valor para mejorar
el reconocimiento \cite{settles2009}.

En este contexto, el presente proyecto propone el dise\~no e
implementaci\'on de \textbf{Cadenza}, una plataforma de digitalizaci\'on
asistida de partituras que integra cuatro componentes: un modelo OMR base
(oemer \cite{oemer2021}) encargado de la transcripci\'on autom\'atica a
MusicXML; un motor de validaci\'on basado en reglas musicales que detecta y
se\~nala errores probables; una interfaz de correcci\'on asistida que
permite al usuario revisar y corregir la transcripci\'on sobre la imagen
original, con reproducci\'on autom\'atica de la partitura digitalizada; y un ciclo
de aprendizaje activo que aprovecha las correcciones
del usuario para priorizar la revisi\'on y mejorar progresivamente el
sistema. La plataforma se orienta a partituras impresas y escritas a mano
actuales, monof\'onicas o de piano simple, excluyendo orquestaciones
completas y manuscritos hist\'oricos, cuyo tratamiento excede el alcance de
este trabajo.

% =====================================================================
% II. PLANTEAMIENTO DEL PROBLEMA
% =====================================================================
\section{Planteamiento del Problema}

En la pr\'actica, la digitalizaci\'on de partituras contin\'ua siendo un
proceso predominantemente reactivo y dependiente del esfuerzo humano.
Aunque el reconocimiento \'optico de m\'usica ha evolucionado desde sistemas
basados en reglas hasta \emph{pipelines} de aprendizaje profundo capaces de
transcribir p\'aginas completas de forma autom\'atica
\cite{omr2020}\cite{calvo2018}\cite{riosvila2026}, su salida a\'un comete
errores frecuentes que impiden usarla directamente \cite{calvo2018}: alturas
o duraciones incorrectas, compases desbalanceados, armaduras inconsistentes
y s\'imbolos omitidos. Por eso, quien digitaliza termina corrigiendo
s\'imbolo por s\'imbolo en un editor convencional, con un costo de tiempo
cercano al de la transcripci\'on manual, o bien acepta resultados sin
verificaci\'on alguna. En el caso de los manuscritos modernos esta situaci\'on se
agrava: la irregularidad del trazo, la inclinaci\'on de las plicas y la
variabilidad geom\'etrica de los s\'imbolos degradan a\'un m\'as la
confiabilidad del reconocimiento autom\'atico \cite{ricordi2024}.

Esta condici\'on evidencia una brecha disciplinar: los modelos de
reconocimiento se estudian predominantemente desde la precisi\'on de su
salida autom\'atica, mientras que su integraci\'on en el trabajo pr\'actico
de un transcriptor ha recibido mucha menos atenci\'on \cite{rizo2023}. El
objeto de estudio de este proyecto es, por tanto, el dise\~no de un sistema
de informaci\'on que resuelva la fragmentaci\'on actual entre tres
dimensiones que hoy operan de manera aislada: (i) los \emph{modelos de
reconocimiento OMR}, evaluados casi exclusivamente mediante m\'etricas de
precisi\'on simb\'olica; (ii) las t\'ecnicas de \emph{validaci\'on musical},
capaces de detectar errores estructurales de m\'etrica, duraciones y
armaduras mediante reglas de teor\'ia musical sin intervenci\'on humana
\cite{cuthbert2010}; y (iii) los \emph{mecanismos de interacci\'on humana},
que permiten al usuario corregir la transcripci\'on sobre la imagen
original, pero cuyas correcciones, en los sistemas actuales, se descartan
sin alimentar con ellas la mejora progresiva del sistema mediante
aprendizaje activo, enfoque pr\'acticamente inexplorado en OMR y con
resultados iniciales incluso contradictorios \cite{amershi2014}%
\cite{settles2009}\cite{mlsp2025}.

A nivel de ingenier\'ia, el problema se concreta en la inexistencia de una
plataforma tecnol\'ogica integrada que articule estas dimensiones. En el
marco experimental de este estudio, las variables se definen de la
siguiente manera: la \emph{Variable Independiente} es la modalidad del
flujo de digitalizaci\'on aplicada a cada entrada, entendida como la
comparaci\'on entre una l\'inea base OMR no asistida y el flujo asistido
con validaci\'on musical y aprendizaje activo, sobre partituras impresas
estructuradas y manuscritas modernas; la \emph{Variable Mediadora} es el
procesamiento anal\'itico de la salida, es decir, la detecci\'on de
anomal\'ias mediante reglas de teor\'ia musical y la selecci\'on activa de
muestras para el reentrenamiento; y las \emph{Variables Dependientes} son
el esfuerzo
de correcci\'on humana, cuantificado en tiempo de edici\'on e
intervenciones manuales por comp\'as, y la fidelidad de la transcripci\'on
final (MusicXML/MIDI), medida mediante la tasa de error simb\'olico
(\emph{Symbol Error Rate}, SER).

En el contexto regional, esta necesidad afecta a archivos, academias y
agrupaciones musicales que conservan colecciones de partituras impresas y
escritas a mano cuya transcripci\'on manual representa costos de tiempo
considerables, y que cuentan \'unicamente con editores convencionales, sin
mecanismos de verificaci\'on autom\'atica ni de mejora continua. La
experiencia reciente de digitalizaci\'on de colecciones musicales,
abordada con metodolog\'ias de anotaci\'on y clasificaci\'on autom\'atica
\cite{ricordi2024}, confirma que el problema es real y abordable desde la
ingenier\'ia. A ello se suma que los conjuntos de datos p\'ublicos con
\emph{ground truth}, como PrIMuS/Camera-PrIMuS \cite{calvo2018}, Sheet
Music Benchmark \cite{smb2026} y MUSCIMA++ \cite{muscima2017}, permiten
evaluar la soluci\'on propuesta de forma cuantitativa y reproducible,
siguiendo los protocolos de la literatura de OMR.

\textbf{Pregunta de Investigaci\'on:} \textquestiondown De qu\'e manera el
dise\~no e implementaci\'on de una plataforma de digitalizaci\'on asistida
de partituras, que integre reconocimiento \'optico de m\'usica, validaci\'on
autom\'atica mediante reglas de teor\'ia musical y un ciclo de aprendizaje
activo a partir de las correcciones del usuario, permite reducir el
esfuerzo de revisi\'on humana y mejorar la fidelidad de la transcripci\'on
digital (MusicXML/MIDI), medida en tasa de error simb\'olico y tiempo de
edici\'on, en partituras impresas y manuscritos modernos de complejidad
monof\'onica o piano simple?

% =====================================================================
% REFERENCIAS — única bibliografía del documento (mapeada a docs/literatura/)
% =====================================================================
\begin{thebibliography}{13}

\bibitem{omr2020}
J. Calvo-Zaragoza, J. Haji\v{c} Jr., and A. Pacha, ``Understanding Optical
Music Recognition,'' \emph{ACM Computing Surveys}, vol.~53, no.~4, pp.~1--35,
2020. % ficha OMR-001

\bibitem{rebelo2012}
A. Rebelo, I. Fujinaga, F. Paszkiewicz, A. R. S. Mar\c{c}al, C. Guedes, and
J. S. Cardoso, ``Optical music recognition: State-of-the-art and open
issues,'' \emph{Int. J. Multimedia Information Retrieval}, vol.~1, no.~3,
pp.~173--190, 2012. % ficha OMR-002

\bibitem{calvo2018}
J. Calvo-Zaragoza and D. Rizo, ``End-to-End Neural Optical Music
Recognition of Monophonic Scores,'' \emph{Applied Sciences}, vol.~8, no.~4,
art.~606, 2018. % ficha DL-001 (introduce PrIMuS)

\bibitem{riosvila2026}
A. R\'ios-Vila, J. Calvo-Zaragoza, and E. Paquet, ``End-to-End Full-Page
Optical Music Recognition for Pianoform Sheet Music,'' \emph{Int. J.
Computer Vision}, 2026, DOI 10.1007/s11263-025-02654-6. % VERIFICAR autores — ficha DL-009

\bibitem{rizo2023}
D. Rizo, J. Calvo-Zaragoza, J. C. Mart\'inez-Sevilla, A. Rosell\'o, and
E. Fuentes-Mart\'inez, ``Design of a music recognition, encoding, and
transcription online tool,'' in \emph{Proc. CMMR 2023}, Tokyo, Japan, 2023;
Springer volume, 2026, DOI 10.1007/978-3-032-02042-0\_24. % ficha HITL-005

\bibitem{cuthbert2010}
M. S. Cuthbert and C. Ariza, ``music21: A toolkit for computer-aided
musicology and symbolic music data,'' in \emph{Proc. 11th Int. Society for
Music Information Retrieval Conf. (ISMIR)}, 2010, pp.~637--642. % ficha VAL-001

\bibitem{amershi2014}
S. Amershi, M. Cakmak, W. B. Knox, and T. Kulesza, ``Power to the people:
The role of humans in interactive machine learning,'' \emph{AI Magazine},
vol.~35, no.~4, pp.~105--120, 2014. % ficha HITL-001

\bibitem{mlsp2025}
``Experimenting Active and Sequential Learning in a Medieval Music
Manuscript,'' in \emph{Proc. IEEE Machine Learning for Signal Processing
Workshop (MLSP)}, 2025, DOI 10.1109/MLSP62443.2025.11204321. % VERIFICAR autores — ficha AL-003

\bibitem{oemer2021}
BreezeWhite, ``oemer: A Python library for end-to-end optical music
recognition,'' GitHub repository, 2021. [Online]. Available:
https://github.com/BreezeWhite/oemer % ficha OMR-003

\bibitem{ricordi2024}
``Optical Music Recognition in Manuscripts from the Ricordi Archive,''
ACM, 2024, DOI 10.1145/3678299.3678324. % VERIFICAR autores — ficha DS-005

\bibitem{smb2026}
J. C. Martinez-Sevilla, J. Cerveto-Serrano, N. Luna, G. Chapman, C. Sapp,
D. Rizo, and J. Calvo-Zaragoza, ``Sheet Music Benchmark (SMB): Standardized
Optical Music Recognition Evaluation,'' 2026, DOI 10.5281/zenodo.17706531.
% ficha FM-005

\bibitem{muscima2017}
J. Haji\v{c} Jr. and P. Pecina, ``The MUSCIMA++ dataset for handwritten
optical music recognition,'' in \emph{Proc. 14th Int. Conf. Document
Analysis and Recognition (ICDAR)}, 2017, pp.~39--46. % ficha DS-009

\bibitem{settles2009}
B. Settles, ``Active Learning Literature Survey,'' Computer Sciences
Technical Report 1648, University of Wisconsin--Madison, 2009. % ficha AL-001

\end{thebibliography}

\end{document}
